\documentclass[12pt]{article}
\usepackage[utf8]{inputenc}
\usepackage[T1]{fontenc}
\usepackage[margin=1in]{geometry}
\usepackage{setspace}
\usepackage{xcolor}
\usepackage[colorlinks=true, urlcolor=blue!60!black, linkcolor=blue!60!black]{hyperref}
\setlength{\parindent}{2em}
\setlength{\parskip}{1.25em}
\renewcommand{\baselinestretch}{1.0}
\begin{document}
\begin{center}
{\LARGE \textit{Youngstown Sheet \& Tube Co. v. Sawyer}, 343 U.S. 579 (1952)\par}
\vspace{0.8em}
{\large Case Brief --- Four-Page Edition\par}
\end{center}
\vspace{1em}

\section{Memory Jogger}

In the landmark ``Steel Seizure Case,'' a 6--3 Supreme Court held that President Truman lacked both statutory and constitutional authority to seize the Nation's strike-bound steel mills during the Korean War, establishing Justice Jackson's celebrated three-zone framework for judging executive power (\href{https://supreme.justia.com/cases/federal/us/343/579/}{343 U.S. 579}).

\section{Detailed Case Facts}

Truman committed troops to Korea in 1950 without a declaration of war and proclaimed a national emergency. In late 1951 the United Steelworkers and the major steel companies deadlocked, and the Wage Stabilization Board's recommendations failed to produce a settlement. When the union called a nationwide strike for April 9, 1952, Truman, believing steel indispensable to weapons and atomic-energy programs, declined to use the Taft-Hartley Act's 80-day cooling-off injunction (a law he had vetoed) or the slower condemnation procedures of the Selective Service and Defense Production Acts.

Instead, on April 8 he issued Executive Order 10340, directing Commerce Secretary Charles Sawyer to seize and operate most of the Nation's steel mills, citing no statute but ``the authority vested in me by the Constitution and laws of the United States, and as President of the United States and Commander in Chief of the armed forces.'' Truman reported the seizure to Congress twice; Congress took no action. The companies obeyed under protest and sued.

\section{Procedural History}

Judge David A. Pine of the District Court for the District of Columbia rejected the claim of unreviewable ``inherent'' emergency power and preliminarily enjoined the seizure. The D.C. Circuit stayed the injunction the same day. The Supreme Court granted certiorari on May 3, heard argument on May 12--13, and affirmed on June 2, 1952.

\section{Judicial Votes}

Affirmed, \textbf{6--3}; every Justice in the majority wrote separately.
\begin{itemize}
\item \textbf{Opinion of the Court:} Justice Black, joined by Justices Frankfurter, Douglas, Jackson, and Burton, each of whom also concurred separately.
\item \textbf{Concurring in the judgment:} Justice Clark.
\item \textbf{Dissent:} Chief Justice Vinson, joined by Justices Reed and Minton.
\end{itemize}

\section{Holding}

The President's power to seize private property must stem from an act of Congress or the Constitution. Neither the Commander-in-Chief, Vesting, nor Take Care Clause authorizes seizing industrial facilities to avert a labor strike, a lawmaking function vested in Congress, especially where Congress has declined to grant seizure power. The case was ripe at the preliminary-injunction stage given irreparable harm. Black's rigid formalism was tempered by the concurrences, and Jackson's framework, technically dicta, later became the governing test in \textit{Dames \& Moore}, \textit{Medell\'{i}n}, and \textit{Zivotofsky}.

\section{Analysis of Opinions}

\textbf{Black (Opinion of the Court).} No statute authorized the seizure; in 1947 Congress had rejected an emergency seizure amendment to Taft-Hartley. The Commander-in-Chief power does not extend ``theater of war'' authority to domestic factories, and the Take Care Clause confines the President to recommending and vetoing laws. The order made policy like a statute and usurped Article I.

\textbf{Frankfurter.} Separation of powers is ``more complicated and flexible.'' Historical practice can supply a gloss on executive power, but the record of sixteen seizure statutes showed Congress withholding general seizure power, and three isolated pre-Pearl Harbor seizures fell far short of \textit{Midwest Oil}.

\textbf{Douglas.} An emergency ``did not create power,'' and separation of powers exists ``not to promote efficiency, but to preclude the exercise of arbitrary power.'' A seizure is a taking, and only Congress can appropriate just compensation.

\textbf{Jackson.} Drawing on his service as Roosevelt's Attorney General, Jackson described three zones: authority ``at its maximum'' when the President acts with congressional authorization; a ``zone of twilight'' when Congress is silent; and power ``at its lowest ebb'' when he acts ``incompatible with the expressed or implied will of Congress.'' Because three statutes covered seizures, Truman was in Zone 3, and the Commander in Chief is not ``Commander in Chief of the country, its industries and its inhabitants.''

\textbf{Burton and Clark.} Burton stressed Congress's deliberate Taft-Hartley procedures. Clark, Truman's former Attorney General, believed in broad emergency power when Congress is silent but required the President to follow the procedures Congress had laid down.

\textbf{Vinson (dissenting).} Citing emergency actions from Washington to Roosevelt, Vinson argued that the Take Care duty included preserving defense programs, that Truman had reported to Congress and would abide its directive, and that the majority adopted a ``messenger-boy concept of the Office.''

\section{Examples: Future Applications}

\begin{enumerate}
\item \textbf{Same-Side (Unconstitutional): Semiconductor Seizure.} The President seizes chip factories to avert a strike, bypassing a labor statute rejecting seizure remedies: Zone 3.
\item \textbf{Same-Side (Unconstitutional): Warrantless Wiretaps.} Domestic wiretapping in violation of FISA places the President's power at its ``lowest ebb.''
\item \textbf{Opposite-Side (Constitutional): IEEPA Asset Freeze.} Freezing foreign assets under statute and long acquiescence falls in Zone 1, as in \textit{Dames \& Moore}.
\item \textbf{Opposite-Side (Constitutional): Recognition.} Under \textit{Zivotofsky}, the President prevails even in Zone 3 when exercising an exclusive recognition power.
\item \textbf{Fence-Sitter: Grid Cybersecurity Directive.} A 30-day order mandating defenses at private utilities amid congressional silence sits in the ``zone of twilight,'' turning on ``the imperatives of events and contemporary imponderables.''
\end{enumerate}

\section{Critique}

\textbf{Formalist and originalist views.} Formalists praise Black's clear defense of Article I's vesting of lawmaking in Congress. Jack Goldsmith and John Manning argue that Black overstated Article II's rigidity and that modern practice partly vindicates Vinson's presidential ``completion power.''

\textbf{Functionalist and doctrinal views.} Functionalists celebrate Jackson's account of the branches' interdependence. Mark Rosen argues that Zone 3 assumes ``categorical congressional supremacy'' without justification, and scholars warn that \textit{Trump v. United States} (2024) risks inverting Jackson's anti-aggrandizement premise.

\section{Key Quotations}

\begin{quote}
``While the Constitution diffuses power the better to secure liberty, it also contemplates that practice will integrate the dispersed powers into a workable government. It enjoins upon its branches separateness but interdependence, autonomy but reciprocity. Presidential powers are not fixed but fluctuate depending upon their disjunction or conjunction with those of Congress.''
\end{quote}
\noindent --- Justice Jackson, concurring, 343 U.S. at 635.

\begin{quote}
``The President's power, if any, to issue the order must stem either from an act of Congress or from the Constitution itself.''
\end{quote}
\noindent --- Justice Black, Opinion of the Court, 343 U.S. at 585.

\begin{quote}
``In the framework of our Constitution, the President's power to see that the laws are faithfully executed refutes the idea that he is to be a lawmaker. The Constitution limits his functions in the lawmaking process to the recommending of laws he thinks wise and the vetoing of laws he thinks bad.''
\end{quote}
\noindent --- Justice Black, Opinion of the Court, 343 U.S. at 587--588.

\end{document}
